\documentclass[12 pt,letterpaper]{article} 

\usepackage[bottom]{footmisc}
\usepackage[utf8]{inputenc}
\usepackage{float}
\usepackage{array}
\usepackage[nottoc]{tocbibind}
\usepackage{caption}
\usepackage{amsmath}
\usepackage[spanish]{babel}
\usepackage{setspace}
\usepackage{fancyhdr}
\usepackage{booktabs}
\usepackage{tabularx}
\usepackage{multirow,hhline}
\usepackage[T1]{fontenc}
\usepackage{lmodern}
\usepackage{amssymb,amsthm}
\usepackage{verbatim} 
\usepackage{comment}
\usepackage{enumitem}
\usepackage{amssymb}
\usepackage{tensor}
\usepackage{lscape}
\usepackage{longtable}

\usepackage[usenames,dvipsnames,svgnames,table]{xcolor}
\usepackage{tabularx}
\usepackage{array}
\usepackage{natbib}
\usepackage{delarray}
\usepackage{dcolumn}
\usepackage{float}
\usepackage{amsmath}
\usepackage{psfrag}
\usepackage{multirow,hhline}
\usepackage{amstext}
\usepackage{pstricks}
\usepackage{pst-plot}
\usepackage{amssymb}
\usepackage{fancyhdr}
\usepackage{dsfont}
\usepackage{transparent}

\usepackage{fontawesome}
\DeclareFontFamily{U}{fontawesomeOne}{}
\DeclareFontShape{U}{fontawesomeOne}{m}{n}
{<-> FontAwesome--fontawesomeone}{}
\DeclareRobustCommand\FAone{\fontencoding{U}\fontfamily{fontawesomeOne}\selectfont} 

\usepackage{framed}
\usepackage[framemethod=tikz]{mdframed}

\usepackage[top=1.5in,bottom=1in,right=1in,left=1in,headheight=65pt]{geometry}

\usepackage[colorlinks=true, linkcolor=magenta]{hyperref}
\hypersetup{
	colorlinks,
	citecolor=Red,
	linkcolor=Cyan,
	urlcolor=Red}


\newif\ifprint
%\printtrue  % Para hacer que se convierta en pauta
% \begin{solution} y \end{solution} para generar respuestas


\ifprint
\newcommand{\tipo}{Pauta}
\newenvironment{solution}
{\begin{mdframed} \textbf{Respuesta:} \ \\}
	{\end{mdframed}}

\else
\newcommand{\tipo}{Enunciado}
\excludecomment{solution}
\fi	

%_____________________________________________________________________________________________________________________



%\usepackage{Sweave}
\begin{document}
%\input{ayudantia_N02_pauta-concordance}

	
	\pagestyle{fancy}
	\fancyhf{}
	\lhead{\center {\transparent{0.4}\includegraphics[width=3cm]{logo_fen.jpg}}}
	\rfoot{Página \thepage}
	\renewcommand{\headrulewidth}{0pt}
	\renewcommand{\footrulewidth}{0pt}	
	
\begin{center}
		
	\bigbreak
	\textbf{ICO 187 ANÁLISIS DE DATOS}\\
	\small{2021}
	\break
	\textbf{CLASE 03: Comprensión de datos numéricos y no numéricos}\\
\end{center}

\begin{flushright}	
		
	Profesor: \href{mailto:sebastianeganasa@santotomas.cl}{Sebastián Egaña}

\end{flushright}

\section{Repaso de la clase pasada}

\begin{itemize}
	\item Referencia relativa: Guarda relación con la columna y fila en donde se encuentran, y son editadas al momento de ser copiadas.
	
	\item Referencia absoluta: La relación con la fila y columna es permanente, considerando la fijación de estas. Se realiza a través de la utilización de nombres, como también a través de la utilización del signo \$.
\end{itemize}
	

\section{Data Númerica}

Se hablo anteriormente de los tipos de datos para Excel, en donde las diferencias para el caso de los números viene dada por las distintas etiquetas que se pueden aplicar a los mismos. Veamos ahora el detalle en el orden de precedencia de las operatorias:

\begin{enumerate}
	\item Evalua los parentesis.
	\item Evalua los rangos.
	\item Evalua las intersecciones (espacio).
	\item Evalua las uniones (;).
	\item Negaciones (-).
	\item Convertir porcentajes.
	\item Exponenciales.
	\item Multiplicación y división, de igual manera (* y /).
	\item Adición y substracción, de igual manera (+ y -).
	\item Operadores de texto (\&).
	\item Comparativos (=, <>, <=, >=).
	
	En el caso de encontrarse elementos en el mismo orden, se evaluan de izquierda a derecha. 
\end{enumerate}

\subsection{Intersecciones y uniones en Excel}

Veamos el siguiente video: \href{https://web.microsoftstream.com/video/a00c02ae-3403-40f1-92ee-a5f5e49843f2}{Enlace acá}


\subsection{Funciones que nos pueden servir}

\begin{itemize}
	\item ESNUMERO: Valor se refiere a un número. \href{https://support.microsoft.com/es-es/office/funciones-es-0f2d7971-6019-40a0-a171-f2d869135665}{Enlace acá}

\end{itemize}

\noindent
{\Huge \faBell} Ejercicio: Desarrolle una función que retorne 1 si la celda contiene un número, y 0 si no lo contine. 

\begin{solution}
	
Considere la siguiente función \\

=SI(ESNUMERO(A1); 1; 2)
	
\end{solution}

\section{Arreglos en Excel}

Desde 2018, Excel incorpora la opción de funciones que se aplican en rangos.

\begin{itemize}
	\item 	Por ejemplo, la función SECUENCIA \href{https://support.microsoft.com/es-es/office/función-secuencia-57467a98-57e0-4817-9f14-2eb78519ca90}{Enlace acá} 
\end{itemize}

\noindent
{\Huge \faBell} Ejercicio: Aplique un arreglo que replique la función SUMAR.SI.

\begin{solution}
Considere la aplicación de las siguientes funciones:\\

=SUMA(SI(\$C\$5:\$C\$11=\$C19;D\$5:D\$11;0))
\end{solution}

\section{Funciones comunes}

\begin{itemize}
	\item SUMA
	\item MIN
	\item MAX
	\item AVERAGE
	\item CONTAR
	\item LARGO
	\item SUMAR.SI
	\item PROMEDIO.SI
	\item CONTAR.SI
	\item DIAS
	\item HOY
\end{itemize}

\section{Funciones matemáticas}

\begin{itemize}
	\item ESNUMERO
	\item ALEATORIO
	\item REDONDEAR
	\item MEDIANA
	\item PI
	\item POTENCIA
	\item RESIDUO
	\item NUMERO.ROMANO
\end{itemize}

\noindent
{\Huge\faComment} Pregunta: La función RESIDUO, por lo general puede ser utilizada para determinar si un número es primo o no, ¿por qué?

\section{Tener en cuenta siempre}

\begin{itemize}
	\item Los valores vacíos, Excel los asume como 0. Ver operaciones. 
	\item Revisar siempre el formato de las fechas.
	\item Revisar siempre el formato númerico; marcador de decimales versus marcador de miles. 
	\item Revisar cuál es el marcador para los parámetros de las funciones. 
\end{itemize}

\section{Desafío \faWarning}

\begin{itemize}
	\item Genere una función que retorno el primer valor no blanco (vacío de un vector). 
	
	Considere la siguiente tabla:
	
	\begin{center}
	\begin{tabular}{|c|c|c|c|}
		\hline
		A & B & C & D \\
		\hline
		100 &  & 100 &  \\
		\hline
		AAA &  &  & AAA \\
		\hline
	\end{tabular}
	\end{center}
	
	En donde la función debería ir en la columna A, identificando el primer valor no blaco y retornarlo en la misma columna. 
	
\end{itemize}

\section{Próxima Clase}

\begin{itemize}
	\item Control Individual de Participación (5\% de ponderación).
\end{itemize}

\section{Aparte}

\begin{itemize}
	\item Ver la disponibilidad de Power Pivot en Excel: \href{https://support.microsoft.com/es-es/office/iniciar-el-complemento-power-pivot-para-excel-a891a66d-36e3-43fc-81e8-fc4798f39ea8}{Enlace acá}	
	
	\item Instalar Power BI en Español: \href{https://powerbi.microsoft.com/es-es/desktop/}{Enlace acá}
	
	\item Crearse una cuenta en R-Studio Cloud: \href{https://rstudio.cloud/}{Enlace acá}
		
\end{itemize}

\section{Fechas relevantes}

\begin{tabular}{|c|l|c|c|}
	\hline
	Unidad & Evaluación & Ponderación & Fecha \\
	\hline
	Unidad I & Evaluación diagnóstica & & 25/03/2021 \\
	\hline
	& Evaluación Individual Participación & (5\%) & 01/04/2021 \\
	\hline
	& Evaluación Grupal & (15\%) & 27/04/2021 \\
	\hline
	& Evaluación Individual - Sumativa I & (30\%) & 11/05/2021 \\
	\hline
	Unidad II & Evaluación Formativa & & 13/05/2021 \\
	\hline
	& Evaluación Individual Participación & (5\%) & 27/05/2021 \\
	\hline
	& Evaluación Grupal & (10\%) & 08/06/2021 \\
	\hline
	& Evaluación Individual - Sumativa II & (15\%) & 17/06/2021 \\
	\hline
	Unidad III & Evaluación Formativa & & 22/06/2021 \\
	\hline
	& Evaluación Individual Participación & (5\%) & 29/06/2021 \\
	\hline
	& Evaluación Individual Sesión I- Sumativa III & (15\%) & 08/07/2021 \\
	\hline
	& Evaluación Individual Sesión II- Sumativa III & (15\%) & 13/07/2021 \\
	\hline
\end{tabular}

\end{document}



